% DISCUSSION - AUTO-UPDATED SECTION
% Last generated: 2026-01-15T07:27:20.548Z

\subsection{Interpretation of Results}

\subsubsection{Participation Dynamics (RQ1)}

The model tests whether participation incentives and delegation policies meaningfully affect turnout and decision quality. Key design implications include:

\begin{itemize}
    \item DAOs with low participation should use lower quorums calibrated to expected turnout
    \item Delegation becomes essential above $\sim$100 members to sustain quorum reach
    \item Dynamic quorum mechanisms may be optimal, adjusting based on recent participation
\end{itemize}

% AUTO-UPDATED: participation findings
Simulation-conditional numeric recommendations are inserted after calibration targets are locked; values here are placeholders until full results are populated.

\subsubsection{Governance Capture Mitigation (RQ2)}

Voting power mitigations create fundamental trade-offs between capture resistance and throughput:

\begin{itemize}
    \item \textbf{Token caps}: Limit whale influence but may reduce large-holder engagement
    \item \textbf{Quadratic thresholds}: More egalitarian, but complex and susceptible to sybil attacks
    \item \textbf{Velocity penalties}: Discourage rapid accumulation but slow legitimate growth
\end{itemize}

No single mitigation dominates; the choice depends on the DAO's threat model and governance goals.

\subsubsection{Proposal Pipeline Design (RQ3)}

The temp-check $\rightarrow$ on-chain pipeline introduces a pre-filtering stage that may reduce on-chain noise at the cost of increased time-to-decision. Key considerations:

\begin{enumerate}
    \item Temp-check stages filter low-quality proposals before they consume on-chain governance bandwidth
    \item Fast-track mechanisms can offset latency for urgent proposals
    \item Abandonment rates may increase if the pipeline creates excessive friction
\end{enumerate}

\subsection{Theoretical Contributions}

This work contributes to theory in several ways:

\begin{enumerate}
    \item \textbf{Participation models}: A framework for quantifying how incentives and delegation affect turnout and decision quality
    \item \textbf{Capture--throughput trade-offs}: Systematic measurement of governance capture mitigation costs
    \item \textbf{Pipeline analysis}: Controlled comparison of proposal filtering mechanisms and their effects on governance efficiency
\end{enumerate}

\subsection{Practical Implications}

For DAO practitioners, this framework suggests:

\begin{enumerate}
    \item \textbf{Calibrate quorums to participation}: Set thresholds based on expected turnout rather than arbitrary values
    \item \textbf{Layer capture mitigations}: Combine caps, quadratic scaling, and velocity penalties rather than relying on one
    \item \textbf{Add temp-check stages}: Pre-filter proposals to reduce on-chain governance load
    \item \textbf{Iterate with data}: Use simulation to test governance changes before on-chain deployment
\end{enumerate}

\noindent
Numeric recommendations should be interpreted as simulation-conditional and comparative rather than universal prescriptions.

\subsection{Connection to Real-World DAOs}

% AUTO-UPDATED: comparison with real DAOs
We include a validation table comparing simulated and empirical turnout; values will be populated after calibration:

\begin{table}[htbp]
\centering
\caption{Comparison with real DAO data}
\begin{tabular}{lccc}
\toprule
DAO & Actual Turnout & Simulated & Diff \\
\midrule
Compound & TBD & TBD & TBD \\
Uniswap & TBD & TBD & TBD \\
Optimism & TBD & TBD & TBD \\
\bottomrule
\end{tabular}
\end{table}

Once populated, this table will indicate how well the agent models reproduce observed participation bands.
